\phantomsection\section*{ВВЕДЕНИЕ}\addcontentsline{toc}{section}{ВВЕДЕНИЕ}

Компьютерная графика является важной составляющей в различных сферах человеческой деятельности, например: киноиндустрия,
разработка видеоигр, архитектурная визуализацая, а также профессиональные симуляторы.

Туман представляет собой сложный оптический эффект, возникающий в результате взаимодействия света с взвешенными частицами воды в атмосфере.
На корректную визуализацию влияют 3 фактора среду: поглощение света, рассеяние луча и затухание света~\cite{fog-effects}. 
Современные методы визуализации
позволяют добиться высокого уровня фотореализма, однако требуют значительных вычислительных ресурсов.

Целью данной работы является разработка программного обеспечения для визуализации сцены из геометрических объектов и неравномерного тумана.

Для достижения поставленной цели необходимо решить следующие задачи:
\begin{itemize}
    \item проанализировать существующие методы моделирования тумана и алгоритмы рендеринга. Провести сравнительный анализ и выбрать подходы для реализации;
    \item разработать алгоритм генерации изображения с поддержкой взаимодействия лучей с туманом;
    \item реализовать алгоритмы генерации изображения;
    \item реализовать модель объемного тумана;
    \item провести исследование влияния параметров алгоритма на производительность.
\end{itemize}

\clearpage
